% !TeX spellcheck = en_GB
% !TeX encoding = UTF-8
% !TeX root = paper.tex
% !TeX TS-program = pdflatex
% !BIB TS-program = biber
%%
%% SiKDD 2026 / IS2026 submission draft.
%% Based on the official IS2026 template (sample-is-2026.tex), trimmed to
%% only what this paper needs. Do not modify margins/typefaces/spacing.
\documentclass[sigconf,natbib=false, language=english]{acmart}

\usepackage[datamodel=acmdatamodel, style=acmnumeric, backend=biber]{biblatex}
\addbibresource{references.bib}

\AtBeginDocument{%
  \providecommand\BibTeX{{%
    Bib\TeX}}}

\setcopyright{rightsretained}
\copyrightyear{2026}
\acmDOI{}
\acmConference[Information Society 2026]{Information Society 2026: 29th international multiconference}{5--9 October 2026}{Ljubljana, Slovenia}
\acmISBN{}
\settopmatter{printccs=false, printacmref=false}

\geometry{a4paper}

\begin{document}

\title[Fusing Supervised and Unsupervised Learning for Hard Drive Failure Prediction]{Fusing Supervised and Unsupervised Learning for Interpretable Hard Drive Failure Prediction on SMART Data}

\author{Luka Petek}
\affiliation{%
  \institution{University of Maribor, Faculty of Electrical Engineering and Computer Science}
  \city{Maribor}
  \country{Slovenia}
}
\email{TODO@student.um.si} % TODO: replace with your actual student email

\author{Vili Podgorelec}
\affiliation{%
  \institution{University of Maribor, Faculty of Electrical Engineering and Computer Science}
  \city{Maribor}
  \country{Slovenia}
}
\email{vili.podgorelec@um.si}

\renewcommand{\shortauthors}{Petek and Podgorelec}

\begin{abstract}
Hard disk drives remain among the most frequently replaced components in storage systems, and their unplanned failure causes data loss, downtime, and increased operational cost. Manufacturer-embedded SMART thresholds detect only 3--10\% of failures at acceptably low false-alarm rates, and most machine-learning approaches to the problem commit to a single learning paradigm -- typically supervised classification, unsupervised anomaly detection, or clustering -- evaluated in isolation. We present a system that instead fuses four independent models trained on the same SMART feature set -- a random forest classifier, an unsupervised autoencoder anomaly detector, a two-stage autoencoder-bottleneck classifier, and a UMAP+HDBSCAN density-based clustering model -- into a single interpretable Health Index Rating (HIR) via a weighted root-mean-square combination. Trained and evaluated on a full year of the Backblaze 2025 dataset (32M+ daily SMART records, 4,414 confirmed failures), the strongest individual component reaches 0.929 ROC-AUC and 89.1\% failure recall, while the componentized HIR score preserves per-model interpretability and remains robust when any single signal is weak or noisy. We discuss the resulting system's deployment as a lightweight, fully offline inference pipeline suitable for data-center fleet monitoring and edge/NAS devices, and outline the trade-offs of ensemble fusion versus single-model deployment.
\end{abstract}

\keywords{hard drive failure prediction, SMART data, ensemble learning, autoencoder, anomaly detection, HDBSCAN clustering, interpretable machine learning, edge AI}

\maketitle

\section{Introduction}
\label{sec:introduction}

Hard disk drives (HDDs) are one of the most commonly replaced hardware components in data centers, and their failure is a leading cause of server-level failure, data loss, service unavailability, and increased operational cost~\cite{TODO:lu2020smarter}. To provide early warning, drive manufacturers embed Self-Monitoring, Analysis, and Reporting Technology (SMART) into virtually every modern drive. However, the threshold-based failure flags built into SMART firmware are deliberately conservative -- manufacturers report failure detection rates of only 3--10\% at an acceptable false-alarm rate of roughly 0.1\% per year~\cite{TODO:murray2005}, leaving a large share of failures undetected until they are catastrophic.

Machine learning applied to raw SMART attribute values, rather than the manufacturer's own threshold flags, has been shown to substantially improve detection rates~\cite{TODO:murray2005,TODO:lu2020smarter}. However, the majority of published approaches commit to a single learning paradigm: a supervised classifier trained on labeled failure/healthy examples, an unsupervised anomaly detector trained only on healthy behaviour, or a clustering-based approach that groups drives by operating regime. Each paradigm has a distinct blind spot -- supervised classifiers are only as good as the (rare, imbalanced) failure examples they were trained on; unsupervised anomaly detectors cannot exploit known failure examples when they exist; clustering approaches provide risk context but no direct failure probability.

\textbf{Contributions.} In this paper we:
\begin{itemize}
    \item combine four independently trained models spanning all three paradigms -- supervised classification, unsupervised anomaly detection, and density-based clustering -- into a single deployable pipeline trained on a full year of the Backblaze 2025 SMART dataset;
    \item propose a weighted root-mean-square fusion, the Health Index Rating (HIR), that preserves the interpretability of each component while being more robust than any single model on its own;
    \item report end-to-end evaluation results (up to 0.929 ROC-AUC, 89.1\% failure recall for the strongest component) and discuss the practical trade-offs of deploying such an ensemble as a lightweight, fully offline system suitable for both data-center fleet monitoring and edge/NAS devices.
\end{itemize}

The remainder of this paper is organized as follows. Section~\ref{sec:related} reviews related work on SMART-based failure prediction. Section~\ref{sec:method} describes the dataset, preprocessing, the four component models, and the HIR fusion formula. Section~\ref{sec:results} reports evaluation results. Section~\ref{sec:discussion} discusses deployment use cases, trade-offs, and limitations, and Section~\ref{sec:conclusion} concludes.

\section{Related Work}
\label{sec:related}
% TODO (next pass): position against Murray et al. 2005 (JMLR), Lu et al. 2020 (FAST, SMARTer),
% Xu et al. 2018 (health-degree, arXiv:1809.04188), the 2023 cost-aware LSTM/imbalance paper,
% and the 2023 TDSC semi-supervised proactive prediction paper. Framing: prior work commits to
% one paradigm; we fuse three into one interpretable score. VERIFY every citation before final
% submission -- do not take these summaries as authoritative without checking the source PDFs.

\section{Data and Methodology}
\label{sec:method}
% TODO (next pass): Backblaze 2025 dataset description, preprocessing/balancing,
% brief subsections for each of the 4 models (RF, AE anomaly, bottleneck AE+classifier,
% UMAP+HDBSCAN), then the HIR weighted-RMS fusion formula and weight rationale.
% Source: README.md and srcML/hir_final.py in the diskFailurePrediction repo (read-only reference).

\section{Results}
\label{sec:results}
% TODO (next pass): performance table (ROC-AUC / Recall / F1 / HIR weight per model),
% 1 key figure, 2-3 sentence interpretation of why fusion helps.

\section{Discussion}
\label{sec:discussion}
% TODO (next pass): deployment use cases (data center fleet monitoring, NAS/edge, offline
% inference), pros (interpretability, no cloud dependency, robustness) vs. cons (current-state
% not time-to-failure, added complexity, no location/environmental features), limitations
% (single dataset/vendor mix, point-in-time snapshot).

\section{Conclusion and Future Work}
\label{sec:conclusion}
% TODO (next pass): summary + future work (time-to-failure horizon, location/environmental
% features, journal extension).

%\begin{acks}
%% TODO: add if there is funding/support to acknowledge; omit section entirely if none.
%\end{acks}

\printbibliography

\end{document}
